% s-187-p187-m003; original PDF 187, 188
\Aphorism{s-187-p187-m003}{278}
Ihr Spiele der Selbstigkeit, wohin seid ihr mir entflohen? Es gab lichte, jugendliche Zeiten, da das reine Ich wie eine saftige Frucht genossen wurde und die Luft wie Champagner schmeckte; noch gar nichts Wirkliches, Erfüllendes war mir geschehen, aber schon die leere Hülse glänzte, als das diamantene Gefäss ungeahnter Zukünfte. Man konnte nur Ich! Ich! sagen, stammeln, lechzen, aber mit diesem hohlen Stirner’schen Ich liess sich Ball spielen; es war eine Seifenblase, aber in paradiesischen Farben. Aus diesem Nichts spann man seine ganze Welt hervor, diese flachste graueste aller Abstractionen gab Hochgebirge her und purpurne Sonnenuntergänge: mit ihr erlebte man Dramen, deren jegliche Person Ich hiess. — Heute ist man älter, realer, von Inhalt beschwerter geworden, und das reine Ich verliert an seinem Wohlgeschmack; nicht mehr wie sonst sieht man stundenlang dem Geflatter eines Schmetterlings zu, der nichts als sich selber trägt. Immer noch lieben wir Beugungen des Lichts und den Schimmer gefärbter Oberflächen, aber dazu wollen wir den beharrenden Edelstein, nicht die schwebende Seifenblase. Immer noch pflücken wir die Frucht, die goldgelbe schwellende, aber wir verstehen nicht mehr aus gemalten Früchten Saft zu saugen. Noch beglückt uns das Ich, das gesättigte, von Welt überfliessende — aber müde sind wir der hundertfachen Spiegelungen eines Ichbegriffs.
\par

